% ===============================
% MAT221 -- Real Analysis
% Riemann Integral
% ===============================

\documentclass[12pt,a4paper]{article}

% ===============================
% Metadata
% ===============================
\newcommand*{\CourseCode}{MAT221}
\newcommand*{\CourseTitle}{Real Analysis}
\newcommand*{\Chapter}{Riemann Integrability}
\newcommand*{\Topics}{Cauchy Criterion, Integration of Continuous Functions}
\newcommand*{\DocDate}{August 17, 2026}

% ===============================
% Header
% ===============================
\input{header}

% ===============================
% Document
% ===============================
\begin{document}

\FrontPage
\Large

% ======================================================
\begin{heading}
Lower Sums
\end{heading}

\begin{definition}[Lower Sum]
Let $f$ be bounded on $[a,b]$, and let
\(
P=\{x_0,x_1,\dots,x_n\}
\)
be a partition of $[a,b]$.

The \textbf{lower sum} of $f$ with respect to $P$ is
\[
L(f,P)
=
\sum_{i=1}^{n}
m_i(x_i-x_{i-1}).
\]
\end{definition}



\clearpage

% ======================================================
\begin{heading}
Upper Sums
\end{heading}

\begin{definition}[Upper Sum]
Let $f$ be bounded on $[a,b]$, and let
\[
P=\{x_0,x_1,\dots,x_n\}
\]
be a partition of $[a,b]$.

The \textbf{upper sum} of $f$ with respect to $P$ is
\[
U(f,P)
=
\sum_{i=1}^{n}
M_i(x_i-x_{i-1}).
\]
\end{definition}

\clearpage

% ======================================================
\begin{heading}
Lower Integral
\end{heading}

\begin{definition}[Lower Integral]
Let $f$ be bounded on $[a,b]$.

The \textbf{lower integral} of $f$ over $[a,b]$ is
\[
\underline{\int_a^b} f
=
\inf
\left\{
L(f,P):
P\text{ is a partition of }[a,b]
\right\}.
\]
\end{definition}

\clearpage

% ======================================================
\begin{heading}
Upper Integral
\end{heading}

\begin{definition}[Upper Integral]
Let $f$ be bounded on $[a,b]$.

The \textbf{upper integral} of $f$ over $[a,b]$ is
\[
\overline{\int_a^b} f
=
\sup
\left\{
U(f,P):
P\text{ is a partition of }[a,b]
\right\}.
\]
\end{definition}

\clearpage

% ======================================================
\begin{heading}
Riemann Integrability
\end{heading}

\begin{definition}[Riemann Integrable Function]
A bounded function
\[
f:[a,b]\to\mathbb{R}
\]
is called \textbf{Riemann integrable} on $[a,b]$ if
\[
\underline{\int_a^b} f
=
\overline{\int_a^b} f.
\]

The common value is called the \textbf{Riemann integral} of $f$ and
is denoted by
\[
\int_a^b f(x)\,dx.
\]
\end{definition}

\clearpage

\begin{problem}
Show directly that the constant function
\[
f(x)=k
\]
is integrable over any closed interval \([a,b]\). What is
\(
\int_a^b f?
\)
\end{problem}


\clearpage

\begin{heading}
    Cauchy Criterion for Riemann Integrability
\end{heading}

\vspace{10pt}

\begin{theorem}[Cauchy Criterion]
A bounded function $f$ on $[a,b]$ is Riemann integrable if and only if for every $\epsilon > 0$, there exists a partition $P$ of $[a,b]$ such that
\[
U(f,P) - L(f,P) < \epsilon.
\]
\end{theorem}

\clearpage

\begin{problem}
Consider \(f(x)=2x+1\) over the interval \([1,3]\). Let \(P\) be the partition consisting of the points
\(
\left\{1,\frac{3}{2},2,3\right\}.
\)

\begin{enumerate}
    \item[(a)] Compute \(L(f,P)\), \(U(f,P)\), and \(U(f,P)-L(f,P)\).

    \item[(b)] What happens to the value of \(U(f,P)-L(f,P)\) when we add the point \(\frac{5}{2}\) to the partition?

    \item[(c)] Find a partition \(P'\) of \([1,3]\) for which
    \[
    U(f,P')-L(f,P')<2.
    \]
\end{enumerate}
\end{problem}

\clearpage

\begin{theorem}[Sequential Criterion]
Prove that a bounded function \(f\) is integrable on \([a,b]\) if and only if there exists a sequence of partitions
    \(
    (P_n)_{n=1}^{\infty}
    \)
    satisfying
    \[
    \lim_{n\to\infty}\bigl[U(f,P_n)-L(f,P_n)\bigr]=0.
    \]
\end{theorem}

\clearpage

\begin{problem}
    Using the previous theorem show that the function
    \[
    f(x)=2x^2+3
\]
    is integrable on \([0,1]\).
\end{problem}

\clearpage


% ======================================================

\begin{problem}
Show that the Dirichlet function $f:\R\to \R$ defined by
\[
f(x)=
\begin{cases}
1,&\text{if }x\in\mathbb{Q},\\
0,&\text{if }x \notin \mathbb{Q},
\end{cases}
\]
is not Riemann integrable on $[0,1]$.
\end{problem}

\clearpage

\begin{heading}
    Continuous Functions are Riemann Integrable
\end{heading}

\vspace{10pt}

\begin{theorem}[Continuity $\implies$ Riemann Integrable]
If a function $f$ is continuous on $[a,b]$, then $f$ is Riemann integrable on $[a,b]$.
\end{theorem}

\clearpage








\EndPage

\end{document}